\chapter*{Abstract}
\addcontentsline{toc}{chapter}{Abstract}
\markboth{Abstract}{Abstract}

Every day, personal belongings --- phones, wallets, identity papers, keys, bags
--- are lost in public spaces, and the mechanisms available to reunite them
with their owners remain largely informal. In Algeria, recovery today depends
on physical lost-property desks with no shared registry, or on social-media
groups where a report is a piece of unstructured text that scrolls out of sight
within hours. The consequence is structural rather than accidental: a lost
report and the corresponding found report may both exist, at the same moment,
and still never meet, because nothing in the system is responsible for
comparing them.

This thesis presents \projectname, a web platform that makes that comparison
automatic. Users report a lost or a found item through a single guided form ---
title, free-text description, category, colour, brand, wilaya, date, and up to
five photographs --- and the platform searches the opposite pool on their
behalf. Matching is \emph{multimodal}. Each report is encoded into a
384-dimensional multilingual sentence embedding, so that a description written
in Arabic, French or English is comparable with one written in either of the
other two; each photograph is encoded into a 512-dimensional CLIP vector, so
that two people photographing the same object from different angles produce
neighbouring points in vector space. Both vectors are stored in PostgreSQL
through the \code{pgvector} extension, which allows approximate nearest
neighbour retrieval and relational filtering --- opposite report type, open
status, plausible date window --- to be expressed in a single query. Retrieved
candidates are then scored by a fusion function that blends the semantic text
similarity with a lexical full-text score, calibrates the raw image similarity,
redistributes weight over whichever modalities are actually present, and adds
small bounded boosts for independent corroborating metadata. The result is a
confidence value, presented to the user as a percentage together with a list of
machine-generated reasons explaining why the suggestion was made.

The system is implemented as a modular monolith: a FastAPI backend exposing a
versioned REST API, an \code{arq} background worker that owns all model
inference so that reporting an item never waits on a neural network, a
PostgreSQL~16 database with \code{pgvector}, Redis as job queue, and a
Next.js frontend serving three fully localised interfaces (English, French and
right-to-left Arabic). Contact details are never exposed by a match alone: a
separate claim workflow requires the claimant to answer verification questions
chosen by the reporter, and identities are exchanged only after the reporter
approves. The prototype covers the complete cycle --- registration,
reporting, embedding, matching, notification, claim, and recovery --- and the
labelled outcome of every confirmed or rejected suggestion is persisted as
training signal for a future calibrated ranking model.

Around this engine, the prototype provides a paid matching tier --- reporting,
searching and claiming stay free, while unlocking a suggestion beyond the first
costs one credit, sold in packs of 200 to 1\,200 dinars through the Chargily
payment gateway --- and an administration console whose every action is written
to an audit log in the same transaction. Measured on its own data, the
multilingual encoder retrieved the right object first in all 72 queries of a
controlled cross-language experiment, against 23 for an English-only model; the
measured distribution of image similarities, with a median of 0.58 between
unrelated photographs, set the calibration of the visual channel; and seven of
the eight suggestions judged by their owners were correct. The whole platform
runs on a CPU in about 1.2\,GB of memory, which allows a public pilot to operate
for about 4\,900 dinars a month, the levies of the auto-entrepreneur status
included.

\vspace{0.8cm}
\noindent\textbf{Keywords:} lost and found, semantic search, sentence
embeddings, CLIP, multimodal matching, vector database, pgvector, approximate
nearest neighbour, FastAPI, Next.js, multilingual information retrieval.
